\emph{Esta sección describe el terreno de la primera generación (1.103 edificios); la versión con la ciudad completa (42.018 huellas de Overture Maps, cauces tallados según la red hidrográfica oficial y unos 852.000 triángulos) se describe en la sección~\ref{sec:res-v4}.} El terreno sin edificios ni vegetación, es decir, el modelo digital del terreno (DTM), es FABDEM V1-2 \citep{hawker2022,fabdem2023}, con resolución de 1 segundo de arco (unos \SI{30}{m}), remuestreado de forma bilineal a una retícula regular de \SI{4}{m} en el sistema UTM, huso 19 sur (datum WGS~84). El remuestreo aumenta la densidad de puntos, pero no mejora la resolución efectiva del modelo de elevación. Esta limitación debe considerarse al interpretar los resultados. Se incorporaron \num{1107} plantas de edificios de OpenStreetMap (OSM; \citealp{osm2026}) como bloques elevados. La altura de cada bloque es la media de $\mathrm{DSM}-\mathrm{DTM}$, donde el modelo digital de superficie (DSM) es Copernicus GLO-30 \citep{copernicus_glo30} y el DTM es FABDEM, cuando la planta mide al menos \SI{900}{m^2} y la diferencia supera \SI{2,5}{m}. Si no se cumple, la altura es el número de pisos de OSM multiplicado por \SI{3}{m}; sin ese dato, se asumen \SI{6}{m}. Para cada bloque se informan la media y el percentil 95 de $\mathrm{DSM}-\mathrm{DTM}$. Se excluyeron cuatro polígonos de OSM etiquetados como edificio que corresponden a recintos (complejos deportivos y un centro comercial con $\mathrm{DSM}-\mathrm{DTM}$ medio inferior a \SI{4}{m}), porque, representados como bloques sólidos en el modelo, habrían desviado artificialmente el flujo. Las calles se rebajaron \SI{0,15}{m}, altura supuesta del cordón, con un ancho según su categoría.
